\exer{Géométrie}
% Banque PT SIA 2013 CLever

\setcounter{exo}{0}
On s'intéresse à un véhicule triporteur permettant de s'incliner en virage.
\begin{center}
\includegraphics[width=\linewidth]{050_01}
\end{center}

On suppose que le mécanisme étudié évolue dans le plan défini par le point O et les directions z0 et x0. Le modèle cinématique adopté est précisé par le schéma cinématique de la figure suivante, sur laquelle sont aussi représentées les données géométriques et les paramètres de mouvements qui seront utilisés dans la question suivante afin de simplifier l'étude.


\begin{center}
\includegraphics[width=\linewidth]{050_02}
\end{center}
  
\question{Déterminer deux équations scalaires reliant alpha1 (avec alpha égal à alpha1 moins alpha10, alpha10 désignant la valeur de alpha1 lorsque l'habitacle n'est pas incliné), theta1 et lambda1. Les directions de projection seront choisies judicieusement. En éliminant theta1, mettre la relation entre alpha1 et lambda1 sous la forme cos(alpha1 + psi) = A/B, puis préciser les expressions de psi, A et B en fonction de a, b, L et lambda1.}
\ifprof
\begin{corrige}

\end{corrige}
\else
\fi


Le tracé de cette relation est laborieux sans moyen numérique. Aussi, il vous est proposé de déterminer la position de certains points de la courbe $\alpha\left( \lambda_1 \right)$ en prenant deux positions d'inclinaison de l'habitacle entre $0^{\circ}$ et $45^{\circ}$. On obtient ainsi sept points pour la plage de variation de $\alpha$ (de $-45^{\circ}$ à $45^{\circ}$). Pour cela, on adopte le paramétrage de la figure suivante en prenant comme origine des angles la position << habitacle non-incliné >>.

\begin{center}
\includegraphics[width=\linewidth]{050_03}
\end{center}
  

\question{Représenter les positions des points $B_1$ et $B_2$ pour les deux positions angulaires choisies. Tracer l'évolution de $\alpha$ en fonction de $\lambda_1$ pour $\alpha$ compris entre $-45^{\circ}$ et $+45^{\circ}$. Est-il possible de décrire cette courbe par une fonction linéaire en prenant comme origine les valeurs des paramètres pour la position <<~habitacle non-incliné~>> (on définit alors le paramètre $\lambda$ tel que : $\lambda = \lambda_1 - \lambda_{10}$) ? Si oui, donner une valeur approximative de sa pente, paramètre noté $R$ pour la suite.}
\ifprof
\begin{corrige}
\end{corrige}
\else
\fi

% BEGIN AUTO CORRIGE INCLUDE
\InclureCorrige{corrige.tex}
% END AUTO CORRIGE INCLUDE